\section{Method}
\label{sec:method}

\subsection{Overview and Representation}
\label{sec:method_overview}

Our pipeline consists of four steps: construct matched interventions,
weight tokens by their intervention responses, average category contrasts,
and evaluate the resulting directions on natural images.
Let a frozen VLM produce the layer-\(\ell\) hidden states at its
visual-token positions,
\begin{equation}
    H^{(\ell)}(I)
    =
    \left[h_1^{(\ell)}(I),\ldots,h_{T(I)}^{(\ell)}(I)\right]^\top
    \in\mathbb{R}^{T(I)\times D},
    \label{eq:visual_representation}
\end{equation}
where \(T(I)\) is the number of visual tokens and \(D\) is the decoder
width. These are the positions occupied by the image representation after
the model's visual-to-language interface, tracked through the decoder;
text and special-token positions are excluded. We estimate one unit
direction \(g_c^{(\ell)}\) per category in a fixed menu
\(\mathcal C\), independently at each analyzed layer. We write
\(C=|\mathcal C|\geq2\) and omit the layer superscript below.

Discovery uses controlled composites; evaluation uses disjoint natural
images. Category annotations and instance masks support intervention
construction and quality control. Spatial annotations do not select the
tokens used for direction estimation or natural-image scoring.
Image-level category labels are accessed to evaluate the frozen scores.
Direction estimation uses visual hidden-state differences and does not
require text embeddings or textual descriptions of the extracted features.

\subsection{Matched Visual Interventions}
\label{sec:matched_interventions}

For a target category \(c\) and competitor \(q_i\neq c\),
we construct triples
\(\mathcal T_i=(I_i^0,I_i^c,I_i^{q_i})\).
The base \(I_i^0\) contains neither category; the other two images contain
an inserted cutout of \(c\) or \(q_i\), respectively. Cutouts are extracted
from annotated instances in the discovery split. Both donor and base
images are disjoint from the natural-image evaluation set.

The two cutouts share a requested paste center and target bounding-box
area, with aspect ratios preserved. Boundary clipping can shift their
realized centers, and foreground-mask areas can differ. The normalized
target area is drawn from a common distribution \(p_A\). We sample sites outside
annotated objects and apply the same alpha blending and local color
harmonization procedure to both insertions. Pixels outside each paste
footprint are unchanged. These controls match the surrounding scene and
insertion procedure, while object shape, appearance, and category-specific
interactions can still differ.

An open-vocabulary detector checks target absence in the base using
threshold \(\tau_{\mathrm{abs}}\). Each composite must contain a detection
of its inserted category with confidence at least
\(\tau_{\mathrm{det}}\) and intersection-over-union at least
\(\tau_{\mathrm{IoU}}\) with the paste footprint. The footprint must also
have nonempty overlap with the input visual-token grid. This spatial
check is confined to intervention quality control.
We retain the observed, gate-passing records and record accepted counts
for each category pair. Filtering and incomplete pair coverage can make
these counts unequal; we do not assume that the retained competitors are
uniformly distributed. Each record contributes to both endpoint categories,
with its contrast oriented toward the category being estimated.
\(N_c\) counts all retained records containing \(c\).
% Report collection-specific p_A, detector prompts, and gate thresholds
% from the discovery configuration; they are not refitted in the ablations.

All three images in a triple have the same dimensions and use identical
preprocessing settings, yielding aligned visual-token positions.
The model, fixed prompt, preprocessing, and hidden-state extraction
convention are also held fixed between discovery and evaluation.

\subsection{Intervention-Response Weighting}
\label{sec:response_weighting}

The paste footprint specifies where the input was edited, but we do not
assume that it specifies the hidden states carrying the category
information. Instead, for each triple, we measure how strongly each
visual token responds to either insertion:
\begin{equation}
    \begin{aligned}
        e_{it}
        &=\tfrac12\bigl(
            \|h_t(I_i^c)-h_t(I_i^0)\|_2\\
        &\hspace{3.5em}
            +\|h_t(I_i^{q_i})-h_t(I_i^0)\|_2
        \bigr),\\
        \alpha_{it}
        &=\frac{e_{it}}{\sum_{j=1}^{T_i}e_{ij}+\epsilon},
    \end{aligned}
    \label{eq:response_weights}
\end{equation}
where \(T_i=T(I_i^0)\) and \(\epsilon=10^{-8}\) stabilizes the denominator.
All visual tokens are eligible. This weighting is symmetric in the two
insertions and assigns greater weight to positions with larger
intervention responses. It provides an estimate of latent support,
although strong responses may also include generic insertion or edit
effects.

We use these same weights for both composites to form the direct contrast
\begin{equation}
    d_i^{c,q_i}
    =
    \sum_{t=1}^{T_i}\alpha_{it}
    \left[h_t(I_i^c)-h_t(I_i^{q_i})\right].
    \label{eq:paired_contrast}
\end{equation}
The base determines the weights; the vector being pooled compares the
two inserted categories. Exchanging the composites preserves the weights
and reverses the contrast. Thus support estimation depends on response
magnitude, while the signed contrast captures how the two categories
differ under matched conditions. A triple with zero response at every
token contributes a zero contrast.

\subsection{Category Contrast Averaging}
\label{sec:direction_estimation}

For each category, we average its contrasts and normalize the result:
\begin{equation}
    \widetilde g_c=\frac{1}{N_c}\sum_{i=1}^{N_c}d_i^{c,q_i},
    \qquad
    g_c=\frac{\widetilde g_c}{\|\widetilde g_c\|_2}.
    \label{eq:direction_average}
\end{equation}
The sign is fixed by subtracting the competitor composite from the target
composite. The resulting direction summarizes the target's contrast
against the sampled concept menu. Normalization is applied after
averaging, so contrast magnitudes contribute to the estimate.
A zero or nonfinite category mean is reported as an extraction failure.
The default estimator gives records equal weight; equal weighting of
observed competitors is a separate ablation
(\cref{sec:competitor_ablation}).

\paragraph{Category-relative interpretation.}
To motivate this estimator, consider the shared-nuisance model
\begin{equation}
    h_t(I_i^z)
    =
    h_t(I_i^0)+\nu_{it}+a_{it}u_z+\xi_{it}^z,
    \quad z\in\{c,q_i\}.
    \label{eq:contrast_decomposition}
\end{equation}
Here \(u_z\) is an idealized category vector, \(\nu_{it}\) collects shared
edit and generic object-presence effects, and \(a_{it}\) is an unknown
support coefficient shared by the two insertions. Residuals
\(\xi_{it}^z\) include departures from these matching assumptions.
Subtraction cancels the shared nuisance in this model, yielding
\begin{equation}
    \begin{aligned}
        d_i^{c,q_i}&=\beta_i(u_c-u_{q_i})+\eta_i,\\
        \beta_i&=\sum_t\alpha_{it}a_{it},\\
        \eta_i&=\sum_t\alpha_{it}
            (\xi_{it}^c-\xi_{it}^{q_i}).
    \end{aligned}
    \label{eq:contrast_identity}
\end{equation}
The support-dependent gain \(\beta_i\) is retained explicitly.
If its conditional expectation is positive and independent of the
competitor, \(\mathbb E[\beta_i\mid c,q_i]=b_c>0\), and the
weighted residual satisfies \(\mathbb E[\eta_i\mid c,q_i]=0\),
a retained competitor distribution \(\pi_c\) gives
\(\mathbb E[\widetilde g_c]=b_c(u_c-\sum_{q\neq c}\pi_c(q)u_q)\).
In the idealized case of uniform coverage over the full menu,
\begin{equation}
    \begin{aligned}
        \mathbb E[\widetilde g_c]
        &=b_c\left(u_c-\frac{1}{C-1}\sum_{q\neq c}u_q\right)\\
        &=b_c\frac{C}{C-1}(u_c-\bar u),
        \qquad \bar u=\frac1C\sum_{j=1}^C u_j.
    \end{aligned}
    \label{eq:centroid_interpretation}
\end{equation}
Under these assumptions, the expected unnormalized estimate aligns with
the category's deviation from the menu mean. Unequal observed pair
frequencies instead define a category-specific reference. This is an
interpretation, not an exact recovery guarantee: response weighting can
introduce competitor-dependent gains and biased residuals, and finite
competitor samples introduce additional variation. We therefore assess
the direction by its transfer to natural images.

\subsection{Natural-Image Readout and Evaluation}
\label{sec:presence_evaluation}

For a held-out natural image \(I\), we compute the cosine alignment
between each visual token and the unit category direction:
\begin{equation}
    r_{tc}(I)
    =
    \frac{h_t(I)^\top g_c}{\max\{\|h_t(I)\|_2,10^{-12}\}}.
    \label{eq:token_alignment}
\end{equation}
The readout uses raw hidden states at the same layer and extraction locus
as discovery. It scores all visual tokens without a mask, bounding box,
or token-level category label.

We aggregate a fixed fraction of the highest-scoring tokens:
\begin{equation}
    \begin{aligned}
        K_I&=\max(1,\lceil\rho T(I)\rceil),\\
        \mathcal P_c(I)
        &=\operatorname{TopKIndices}_t(r_{tc}(I),K_I),\\
        s_c(I)&=\frac1{K_I}\sum_{t\in\mathcal P_c(I)}r_{tc}(I).
    \end{aligned}
    \label{eq:presence_score}
\end{equation}
We use \(\rho=0.05\) as the default pooling fraction, fixed without
evaluation-set tuning. Equal values at the cutoff do not change the
pooled score. Top-fraction
pooling allows a subset of tokens to provide evidence for category
presence and scales the pooling budget with the image's token count.
It does not guarantee invariance to resolution or object size.

Let \(\mathcal E\) be the natural-image evaluation set and
\(y_c(I)\in\{0,1\}\) indicate whether category \(c\) is annotated as
present. We evaluate the frozen scores using
\begin{equation}
    \begin{aligned}
        \operatorname{AUC}_c
        &=\operatorname{ROC\!-\!AUC}
            \bigl(\{y_c(I)\}_{I\in\mathcal E},
                  \{s_c(I)\}_{I\in\mathcal E}\bigr),\\
        \operatorname{AUC}_{\mathrm{macro}}
        &=\frac1C\sum_{c\in\mathcal C}\operatorname{AUC}_c.
    \end{aligned}
    \label{eq:presence_auc}
\end{equation}
The headline statistic is the raw macro AUC. It requires a valid
direction and both positive and negative evaluation examples for every
category in the fixed menu; invalid categories are reported explicitly
and are not silently dropped from the average.

As a control, we draw independent Gaussian vectors
\(v_c^{(b)}\sim\mathcal N(0,I_D)\), normalize them to unit length, and
substitute them for \(g_c\) in the identical readout. We use 25 repetitions
and report the mean and standard deviation of their macro AUCs separately.
This control measures the association obtained from random directions
under the same token geometry and pooling rule. Presence AUC evaluates
the transfer of an intervention-derived direction; it does not establish
causal use of that direction in the model's generated answers.
